\scanpage{51}
\mspage{\deleted{\unsure{6}}7}
of an already oriented structure if $a$ is a free edge
(corresponding \unsure{to|:} a $C\in\pi_0(\partial X)$), or that one is in possession
\unsure{\inserted{of}|of} one of the two orientations $\omega(a)$ ---
in the case where $a$ is an ordinary edge, indicates that one has in hand,
for each $a\in A$, a combinatorial polygon $P_a$, whose set of vertices
is $R_a$, and the set of edges can be described \deleted{\unsure{also|as}}
in a canonical way in terms of $R_a,\omega_a$ and
$\omega_a\longrightarrow\operatorname{Circ}(R_a)$.
It may be more convenient to give it from the outset in the structure,
where the family of $(R_a)$ would be replaced by a
\deleted{which would be \unsure{in a contour|the ones a contour}}
\uline{combinatorial contour}, i.e. a set of frames
$\vec{\underline R}$ (rather ``rigidification frames''), with an action of
\[
 \mathbb D_\infty=\{\sigma_0,\sigma_1\mid\sigma_0^2=\sigma_1^2=1\}
\]
on it, and a map
\[
 \vec{\underline R}\longrightarrow\vec A
\]
compatible with the homomorphism
\[
 \mathbb D_\infty\xrightarrow{\varepsilon}\{\pm1\}
\]
on the groups of operators, but in addition a subset
\[
 \vec{\underline R}_l^+\subset\vec{\underline R}\mid\vec A_l
\]
\margin{$(\vec A_l=\text{free edges, corresponding to the gluing circles, i.e. the }\in\pi_0(\unsure{\partial}X))$}
defining an orientation of the corresponding contours.
\mspage{50}
